An orthogonal basis exists by \claimref{std-algebra.sylvester-inertia}. Write each $v_j = \sum_i a_{ji} e_i$ and expand $[v_1, \dots, v_k]$ multilinearly: it is $\sum_{i_1, \dots, i_k} a_{1 i_1} \cdots a_{k i_k} [e_{i_1}, \dots, e_{i_k}]$. A term with a repeated index vanishes, since the antisymmetrisation of a word containing $e_i$ twice cancels in pairs; a term with distinct indices equals $\pm e_I$ for $I = \{i_1, \dots, i_k\}$, because distinct orthogonal basis vectors anticommute (\claimref{std-algebra.orthogonal-anticommute}) and so every reordering is the product times its sign. Hence $\Cl^{k} \subseteq \operatorname{span}\{e_I : |I| = k\}$, and conversely $e_I = [e_{i_1}, \dots, e_{i_k}] \in \Cl^{k}$. The blades form a basis (\claimref{std-algebra.blade-basis}), so the sum over $k$ is direct and $\dim \Cl^{k} = \binom{n}{k}$. Basis independence holds because $\Cl^{k}$ was defined without a basis. The last two formulas: $uv = \tfrac12(uv + vu) + \tfrac12(uv - vu) = B(u, v) + [u, v]$.
